\subsection{Equinumerous Sets}

\newpage
\begin{tcolorbox}[title=Problem 1, breakable]
    Show that the following sets are denumerable.
    \begin{enumerate}
        \item $\mathbb{N}$.
        \item The set of all even integers.
    \end{enumerate}
\end{tcolorbox}

\begin{proof}
    Consider the mapping $f : \mathbb{Z}^+ \longrightarrow \mathbb{N}$ defined 
    \[f(x) = x - 1.\]
    Let $a, b \in \mathbb{Z+}$ and suppose 
    $f(a) = f(b)$. Then 
    \[f(a) = f(b) = a - 1 = b - 1 \iff a = b.\]
    Let $c \in \mathbb{N}$ and notice $c + 1 \in \mathbb{Z}^+$, then 
    \[f(c + 1) = (c + 1) - 1 = c.\]
    Thus $\mathbb{N}$ is denumerable.
\end{proof}

\begin{proof}
    Consider the mapping $f : \mathbb{Z}^+ \longrightarrow 2\mathbb{Z}$ defined 
    \[f(x) = \begin{cases} x - 1 & x \text{ odd} \\ -x & x \text{ even} \end{cases}\]
    Let $a, b \in \mathbb{Z}^+$ and suppose $f(a) = f(b)$. If both are odd, 
        $a - 1 = b - 1 \iff a = b$. If both are even, $-a = -b \iff a = b$. 
    If they have opposite parity, one output is $\geq 0$ and the other 
        is $< 0$, so $f(a) = f(b)$ is impossible. 
    Thus $f$ is injective.
    Let $c \in 2\mathbb{Z}$. 
    If $c \geq 0$ then $x = c + 1$ 
        is odd and $f(c + 1) = c$. If $c < 0$ then $x = -c$ is even and $f(-c) = c$. 
    Thus $f$ is surjective, and $2\mathbb{Z}$ is denumerable.
\end{proof}